\documentclass[troisieme]{fiche}
\metadonnees{11}{La tache de sang}{Bloc C — Décrire et comparer}

\begin{document}
\entete

\objectifs{%
\textbf{Langue} : le comparatif ; \emph{as\ldots as} et \emph{not as\ldots as}.\par
\textbf{Texte} : Oscar Wilde, \og The Canterville Ghost \fg{} (1887), ch. \textsc{i}.\par
\textbf{Durée} : 1 h — exercices 20 min, texte et questions 25 min, version 15 min.}

%% ===================================================================
\exercice[12 min]{Comparer}

\rappel{\textbf{Rappel.} Adjectif court : \emph{-er\ldots than}. Adjectif long :
\emph{more\ldots than}. Égalité : \emph{as\ldots as}.}

\consigne{Complétez avec l'adjectif ou l'adverbe entre parenthèses.}

\begin{anglais}
\begin{enumerate}[label=\arabic*.,itemsep=2.5mm,leftmargin=*]
  \item Canterville Chase is much \trou{older} than the Otis house. \emph{(old)}
  \item The library was \trou{darker} than the Tudor hall. \emph{(dark)}
  \item Mrs Umney's voice became \trou{more mysterious} at every word. \emph{(mysterious)}
  \item The twins are \trou{noisier} than their sister. \emph{(noisy)}
  \item Washington dances \trou{better} than anybody in London. \emph{(well)}
  \item The second stain was \trou{worse} than the first. \emph{(bad)}
  \item The new broom is \trou{as good as} the old one. \emph{(good, égalité)}
  \item The library was \trou{as dusty as} an old church. \emph{(dusty, égalité)}
  \item After the journey the children were \trou{as weary as} their parents.
    \emph{(weary, égalité)}
  \item Mr Otis was \trou{not so frightened as} his housekeeper.
    \emph{(frightened, égalité niée)}
\end{enumerate}
\end{anglais}

\note[Former le comparatif]{Une syllabe : \emph{-er} (1, 2). Deux syllabes en \emph{-y} :
\emph{-ier}, le \emph{y} devient \emph{i} (4, 8). Trois syllabes et plus : \emph{more} (3).
Irréguliers : \emph{good} et \emph{well} donnent \emph{better}, \emph{bad} donne \emph{worse}
(5, 6).}

\note[L'égalité niée]{\emph{not as\ldots as} et \emph{not so\ldots as} se disent tous les
deux ; la seconde forme est un peu plus soutenue. À la forme affirmative, en revanche, on
n'emploie que \emph{as\ldots as}.}

%% ===================================================================
\exercice[8 min]{Renforcer une comparaison}

\rappel{\textbf{Rappel.} \emph{Much} renforce un comparatif (\emph{much bigger}) ;
\emph{very} ne le fait jamais. \og De plus en plus \fg{} : \emph{-er and -er}.}

\consigne{Traduisez en anglais.}

\begin{anglais}
\begin{enumerate}[label=\arabic*.,itemsep=3mm,leftmargin=*]
  \item Il fait de plus en plus froid.
    \lignes{1}
    \corr{It is getting colder and colder.}
  \item Cette maison est beaucoup plus grande que la nôtre.
    \lignes{1}
    \corr{This house is much bigger than ours.}
  \item Plus il criait, moins on l'écoutait.
    \lignes{1}
    \corr{The more he shouted, the less they listened to him.}
  \item Elle court deux fois plus vite que moi.
    \lignes{1}
    \corr{She runs twice as fast as I do.}
  \item Ce n'est pas aussi simple que tu le crois.
    \lignes{1}
    \corr{It is not as simple as you think.}
  \item Il est moins riche que son oncle.
    \lignes{1}
    \corr{He is not as rich as his uncle. \emph{ou} He is less rich than his uncle.}
\end{enumerate}
\end{anglais}

\note[Trois tournures à retenir]{\og De plus en plus \fg{} : l'adjectif répété,
\emph{colder and colder}, \emph{more and more difficult}. \og Plus\ldots, moins\ldots \fg{} :
\emph{the more\ldots, the less\ldots}, avec l'article \emph{the} devant chaque terme.
\og Deux fois plus \fg{} : \emph{twice as\ldots as}, et non \emph{two times}.}

\note[\emph{Less} ou \emph{not as}]{L'anglais emploie \emph{less} beaucoup moins souvent que
le français n'emploie \og moins \fg. Devant un adjectif court, \emph{not as rich as} est plus
naturel que \emph{less rich than}.}

%% ===================================================================
\newpage
\section{Texte}

\noindent\small\textit{Une famille américaine, les Otis, vient d'acheter un manoir anglais
que tout le monde dit hanté. Ils arrivent un soir de juillet.}\normalsize

\begin{texte}
Standing on the steps to receive them was an old woman, neatly dressed in black silk, with a
white cap and \gl[an apron]{apron}{un tablier}. This was Mrs. Umney, the
\gl[a housekeeper]{housekeeper}{une gouvernante}, whom Mrs. Otis, at Lady Canterville's
earnest request, had consented to keep in her former position. She made them each a low
\gl[a curtsey]{curtsey}{une révérence} as they \gl[to alight]{alighted}{descendre de
voiture}, and said in a \gl[quaint]{quaint}{désuet, pittoresque}, old-fashioned manner, ``I
bid you welcome to Canterville Chase.'' Following her, they passed through the fine Tudor
hall into the library, a long, low room, \gl[to panel]{panelled}{lambrissé} in black oak, at
the end of which was a large stained glass window. Here they found tea laid out for them,
and, after taking off their \gl[a wrap]{wraps}{un châle, un manteau}, they sat down and began
to look round, while Mrs. Umney \gl[to wait on]{waited on}{servir} them.

Suddenly Mrs. Otis caught sight of a \gl[dull]{dull}{terne} red \gl[a stain]{stain}{une tache}
on the floor just by the \gl[a fireplace]{fireplace}{une cheminée}, and, quite unconscious of
what it really signified, said to Mrs. Umney, ``I am afraid
something has been \gl[to spill, spilt]{spilt}{renverser} there.''

``Yes, madam,'' replied the old housekeeper in a low voice, ``blood has been spilt on that
spot.''

``How \gl[horrid]{horrid}{affreux}!'' cried Mrs. Otis; ``I don't at all care for blood-stains
in a sitting-room. It must be removed at once.''

The old woman smiled, and answered in the same low, mysterious voice, ``It is the blood of
Lady Eleanore de Canterville, who was murdered on that very spot by her own husband, Sir
Simon de Canterville, in 1575. Sir Simon survived her nine years,
and disappeared suddenly under very mysterious circumstances. His body has never been
discovered, but his \gl[guilty]{guilty}{coupable} spirit still \gl[to haunt]{haunts}{hanter}
the Chase. The blood-stain has been much admired by tourists and others, and cannot be
removed.''

``That is all \gl[nonsense]{nonsense}{des sottises}!'' cried Washington Otis; ``Pinkerton's
Champion Stain Remover and Paragon Detergent will clean it up
in no time,'' and before the terrified housekeeper could
\gl[to interfere]{interfere}{intervenir}, he had fallen upon his knees, and was rapidly
\gl[to scour]{scouring}{récurer} the floor with a small stick of what looked like a black
cosmetic. In a few moments no trace of the blood-stain could be seen.

``I knew Pinkerton would do it,'' he exclaimed, triumphantly, as he looked round at his
admiring family; but no sooner had he said these words than a terrible flash of
\gl[lightning]{lightning}{la foudre, les éclairs} lit up the sombre room,
a fearful \gl[a peal of thunder]{peal}{un coup (de tonnerre)} of thunder made them all start
to their feet, and Mrs. Umney \gl[to faint]{fainted}{s'évanouir}.
\end{texte}
\source{Oscar Wilde, \og The Canterville Ghost \fg, dans \emph{Lord Arthur Savile's Crime and
Other Stories}, Londres, Osgood, 1891, ch.~\textsc{i}.}

\partiel[15 min]{Questions}
\consigne{Répondez aux deux premières questions en français, aux deux dernières en anglais.}

\begin{enumerate}[label=\arabic*.,itemsep=2.5mm,leftmargin=*]
  \item Que remarque Mrs Otis dans la bibliothèque, et quelle explication la gouvernante lui
    en donne-t-elle ?
    \lignes{3}
    \corr{Une tache rouge terne sur le plancher, près de la cheminée. C'est, dit la
    gouvernante, le sang de Lady Eleanore de Canterville, assassinée à cet endroit même par
    son mari en 1575 ; l'esprit coupable de celui-ci hante toujours la maison, et la tache ne
    peut pas être enlevée.}
  \item Comment réagit Washington Otis ? Qu'est-ce que sa réaction montre de la famille ?
    \lignes{3}
    \corr{Il traite l'histoire de sottises, se jette à genoux et frotte la tache avec un
    détachant de marque, qui la fait disparaître aussitôt. La famille répond à une légende de
    trois siècles par un produit du commerce : elle ne croit qu'à ce qui s'achète et
    s'explique.}
  \item \en{Mrs Umney speaks ``in a low, mysterious voice''. Find two other details that show
    she treats the stain as something valuable.}
    \lignes{3}
    \corr{She says the stain \emph{has been much admired by tourists and others}, as if it
    were a work of art in a museum, and she gives the full name and the exact date, like a
    guide. She also smiles before answering: she is enjoying the effect.}
  \item \en{Find the sentence that begins with \emph{no sooner}. Which of the two events
    happens first?}
    \lignes{2}
    \corr{\emph{No sooner had he said these words than a terrible flash of lightning lit up
    the sombre room.} Washington speaks first, mais à peine : la tournure dit que le second
    événement suit immédiatement le premier. Après \emph{no sooner} placé en tête, l'auxiliaire
    passe devant le sujet — \emph{had he said}.}
\end{enumerate}

\note[Une comparaison cachée]{\emph{No sooner\ldots than} est un comparatif : \emph{soon}
donne \emph{sooner}, et \emph{than} introduit le second terme. Le français dit \og à peine\ldots
que \fg.}

%% ===================================================================
\section{Version}
\consigne{Traduisez, de \emph{The old woman smiled} à \emph{could be seen}.}

\lignes{14}

\corr{\textbf{Proposition de traduction.} La vieille femme sourit et répondit de la même voix
basse et mystérieuse : \og C'est le sang de Lady Eleanore de Canterville, qui fut assassinée à
cet endroit même par son propre mari, Sir Simon de Canterville, en 1575. Sir Simon lui survécut
neuf ans, puis disparut soudain dans des circonstances très mystérieuses. On n'a jamais
retrouvé son corps, mais son esprit coupable hante encore le domaine. La tache de sang a été
fort admirée par les touristes et par d'autres, et elle ne peut pas être enlevée. \fg{} \og
Tout cela n'est que sottises ! \fg{} s'écria Washington Otis. \og Le Détachant Champion et
Détergent Paragon de chez Pinkerton en viendra à bout en un rien de temps \fg ; et, avant que
la gouvernante terrifiée pût intervenir, il était tombé à genoux et frottait vivement le
plancher avec un petit bâtonnet qui avait l'air d'un fard noir. En quelques instants, il ne
restait plus trace de la tache de sang.}

\note[Choix de traduction 1]{\emph{His body has never been discovered} : present perfect
passif. Le français passe par \og on \fg{} et le passé composé actif — \og on n'a jamais
retrouvé son corps \fg{} — plus naturel que le passif.}

\note[Choix de traduction 2]{\emph{before the terrified housekeeper could interfere} :
\emph{could} marque ici le temps qui a manqué, non la capacité. \og Avant que \fg{} appelle le
subjonctif en français.}

\note[Choix de traduction 3]{\emph{a small stick of what looked like a black cosmetic} :
\emph{what} = \og ce qui \fg. Le narrateur ne nomme pas l'objet, il dit seulement à quoi il
ressemblait ; la traduction doit garder cette prudence.}

%% ===================================================================
\partiel{Lexique à mémoriser}
\begin{multicols}{2}\small
\begin{itemize}[label=--,itemsep=0.8mm,leftmargin=*]
  \item \textbf{a housekeeper} — une gouvernante
  \item \textbf{a stain} — une tache
  \item \textbf{blood} — le sang ; \textit{adjectif} bloody
  \item \textbf{to spill} — renverser ; spilt, spilt
  \item \textbf{a fireplace} — une cheminée
  \item \textbf{to remove} — enlever, ôter
  \item \textbf{nonsense} — des sottises
  \item \textbf{a knee} — un genou
  \item \textbf{lightning} — la foudre, les éclairs
  \item \textbf{thunder} — le tonnerre
  \item \textbf{to faint} — s'évanouir
  \item \textbf{dull} — terne ; \textit{et} ennuyeux
\end{itemize}
\end{multicols}

\end{document}
